\chapter{Trust boundary}\label{chap:trust}

\section*{At a glance}
\begin{itemize}
\item \lead{Goal} State what is proved about the eraser, what the eraser is made of, and what a
  result built on lean4lean inherits.
\item \lead{Main results} None: \hl{no theorem about the eraser exists}.
\item \lead{Status} \stShipping{} the shipping code, described by definition nodes;
  \stInherited{} the sorry sources and axioms of lean4lean, tabulated, reached by no node.
\item \lead{Lean files} \code{Basic.lean}, \code{Erasure.lean} and \code{Printing.lean} under
  \code{LeanToLambdaBox/}.
\end{itemize}

\section{What is proved}\label{sec:trust-proved}

\begin{itemize}
\item \lead{Theorems} None. The Lean library of the repository is the shipping code; no
  declaration states a property of the eraser.
\item \lead{Hypotheses} None, since there is no statement.
\item \lead{Divergences} None: the register is empty (Chapter~\ref{chap:divergences}).
\item \lead{Tests} The shipping regression tests (\code{scripts/regress.sh}) compare the eraser's
  outputs on fixed inputs with committed expected files.
\end{itemize}

\section{The shipping code}\label{sec:trust-shipping}

The eraser is Lean meta-code: it runs in the elaborator's monads and asks the elaborator for types.
Its recursive traversal consists of \code{partial} definitions, which the kernel cannot unfold.

\begin{definition}[Erased terms]
  \label{def:LBTerm}
  \lean{LBTerm}
  \leanok
  \srcloc{LeanToLambdaBox/Basic.lean}{87}
  \stShipping{} \lead{In short} The $\lambda_\Box$ terms the eraser builds: MetaRocq's
  \code{EAst.term} (\code{erasure/theories/EAst.v}) with \hl{free variables added}.
  \begin{itemize}
  \item \code{box}, \code{bvar}, \code{lambda}, \code{letIn}, \code{app}, \code{const},
        \code{construct}, \code{case}, \code{proj}, \code{fix}, \code{prim}: the MetaRocq
        constructors \code{tBox}, \code{tRel}, \code{tLambda}, \code{tLetIn}, \code{tApp},
        \code{tConst}, \code{tConstruct}, \code{tCase}, \code{tProj}, \code{tFix}, \code{tPrim};
  \item \code{fvar}: a Lean free variable, used during erasure and abstracted to \code{bvar}
        before a binder is built;
  \item absent: \code{tVar}, \code{tEvar}, \code{tCoFix}, \code{tLazy}, \code{tForce}.
  \end{itemize}
\end{definition}

\begin{definition}[Global declarations]
  \label{def:GlobalDecl}
  \lean{GlobalDecl, ConstantBody, MutualInductiveBody, OneInductiveBody, ConstructorBody,
    ProjectionBody}
  \leanok
  \uses{def:LBTerm}
  \srcloc{LeanToLambdaBox/Basic.lean}{171}
  \stShipping{} \lead{In short} The declarations of an erased environment.
  \begin{itemize}
  \item \code{constantDecl}: a constant with an optional body (none for an axiom);
  \item \code{inductiveDecl}: a block of mutual inductive types: parameter count, and per type
        its name, propositional flag, allowed eliminations, constructors and projections.
  \end{itemize}
\end{definition}

\begin{definition}[Erased programs]
  \label{def:ASTType}
  \lean{ASTType, Program}
  \leanok
  \uses{def:LBTerm, def:GlobalDecl}
  \srcloc{LeanToLambdaBox/Basic.lean}{179}
  \stShipping{} \lead{In short} What \code{\#erase} emits: \code{untyped} of a list of named global
  declarations and an optional main term, the \code{Untyped} program of Peregrine's format.
\end{definition}

\begin{definition}[Erasure configuration]
  \label{def:Erasure-ErasureConfig}
  \lean{Erasure.ErasureConfig, Erasure.Config.Extern, Erasure.Config.Nat}
  \leanok
  \srcloc{LeanToLambdaBox/Erasure.lean}{68}
  \stShipping{} \lead{In short} The options of \code{\#erase ... config}.
  \begin{itemize}
  \item \code{extern}: an \code{@[extern]} constant becomes an axiom (\code{preferAxiom}, the
        default) or keeps its Lean definition (\code{preferLogical});
  \item \code{nat}: \code{Nat} literals become primitive integers (\code{machine}, the default) or
        constructor terms (\code{peano});
  \item \code{csimp}, \code{remove\_irrel\_constr\_args},
        \code{auto\_inline\_typeclass\_dispatch}: Boolean options.
  \end{itemize}
\end{definition}

\begin{definition}[The erasure monad]
  \label{def:Erasure-EraseM}
  \lean{Erasure.EraseM, Erasure.ErasureState, Erasure.ErasureContext, Erasure.run}
  \leanok
  \uses{def:GlobalDecl, def:Erasure-ErasureConfig}
  \srcloc{LeanToLambdaBox/Erasure.lean}{168}
  \stShipping{} \lead{In short} $\mathrm{EraseM} = $ \code{StateT ErasureState (ReaderT
  ErasureContext CoreM)}: \hl{Lean's core monad}, with a state and a context.
  \begin{itemize}
  \item state: the inductive types and constants registered so far, the emitted declarations, the
        constants to inline and their inlined sizes;
  \item context: a Lean local context, the fixpoint variables in scope, the configuration.
  \end{itemize}
\end{definition}

\begin{definition}[Erasability oracle]
  \label{def:Erasure-isErasable}
  \lean{Erasure.isErasable}
  \leanok
  \srcloc{LeanToLambdaBox/Erasure.lean}{183}
  \stShipping{} \lead{In short} A term is erased to $\Box$ iff its type is a proposition or a
  type former, as \hl{decided by the elaborator}.
  \begin{itemize}
  \item \lead{Lean} a \code{MetaM} function calling \code{Meta.inferType}, \code{Meta.isProp} and
        \code{Meta.isTypeFormerType};
  \item \lead{Caveat} these belong to Lean's elaborator (\code{Lean.Meta}), a different program
        from the kernel that lean4lean models.
  \end{itemize}
\end{definition}

\begin{definition}[Preprocessing]
  \label{def:Erasure-prepare_erasure}
  \lean{Erasure.prepare_erasure}
  \leanok
  \uses{def:Erasure-EraseM}
  \srcloc{LeanToLambdaBox/Erasure.lean}{438}
  \stShipping{} \lead{In short} Rewrites the input term with passes of Lean's compiler before
  erasure.
  \begin{itemize}
  \item \lead{Steps} unsafe-recursion names replaced, \code{macroInline}, \code{inlineMatchers},
        \code{macroInline} again, then (option \code{csimp}) the \code{@[csimp]} replacements;
  \item \lead{Caveat} the result may be ill-typed when a type depends on a function that
        \code{csimp} replaces (its docstring).
  \end{itemize}
\end{definition}

\begin{definition}[The eraser]
  \label{def:Erasure-erase}
  \lean{Erasure.erase, Erasure.erase.visitExpr}
  \leanok
  \uses{def:LBTerm, def:ASTType, def:Erasure-ErasureConfig, def:Erasure-EraseM,
    def:Erasure-isErasable, def:Erasure-prepare_erasure}
  \srcloc{LeanToLambdaBox/Erasure.lean}{461}
  \stShipping{} \lead{In short} Erases a term and its dependencies to a program; the traversal is
  \hl{partial, hence opaque to proofs}.
  \begin{itemize}
  \item \code{erase e config}: runs \code{prepare\_erasure} then \code{visitExpr} in
        \code{EraseM}; returns the program (the emitted declarations and the main term) and the
        constants to inline;
  \item \code{visitExpr}: an erasable term becomes $\Box$; otherwise one case per \code{Expr}
        constructor, through the \code{where} functions of \code{erase};
  \item \lead{Lean} \code{erase} is a \code{CoreM} definition; \code{visitExpr} and its siblings
        are \code{partial def}s (Section~\ref{sec:trust-census}).
  \end{itemize}
\end{definition}

\begin{definition}[Printing]
  \label{def:Serialize}
  \lean{Serialize, sexpr, LBTerm.to_sexpr}
  \leanok
  \uses{def:LBTerm, def:ASTType}
  \srcloc{LeanToLambdaBox/Printing.lean}{15}
  \stShipping{} \lead{In short} The S-expression printer of programs, in Peregrine's \code{.ast}
  format.
  \begin{itemize}
  \item \code{Serialize}: a class with one method, \code{to\_sexpr}, to the S-expressions
        \code{sexpr};
  \item \code{LBTerm.to\_sexpr}: a \code{partial def}; it reaches \code{unreachable!} on
        \code{fvar}.
  \end{itemize}
\end{definition}

\begin{definition}[The command \texorpdfstring{\code{\#erase}}{\#erase}]
  \label{def:Erasure-eraseElab}
  \lean{Erasure.eraseElab}
  \leanok
  \uses{def:Erasure-erase, def:Serialize, def:Erasure-ErasureConfig}
  \srcloc{LeanToLambdaBox/Erasure.lean}{869}
  \stShipping{} \lead{In short} Elaborates the term, erases it, and writes the program, the
  constants to inline and the \code{.mli} signature: to files when \code{to} and \code{mli} are
  given, to the message log otherwise.
  \begin{itemize}
  \item the \code{config} term is evaluated with \code{unsafe Elab.Term.evalTerm};
  \item the \code{.mli} signature comes from the Lean type of the term (\code{gen\_mli}).
  \end{itemize}
\end{definition}

\subsection{Partial, opaque, unsafe and monadic declarations}\label{sec:trust-census}

The table is generated from the Lean environment by \code{blueprint/scripts/audit.py}.

% Generated by blueprint/scripts/audit.py --update from the Lean environment. Do not edit.
\begin{itemize}
\item \lead{Rows} the declarations of the modules \code{LeanToLambdaBox*} that are partial, opaque or unsafe, or whose result type is a monad of the elaborator or of IO: 10 def, 1 opaque, implemented\_\allowbreak{}by, 11 partial def, 1 unsafe def.
\item \lead{Partial def} an opaque constant for the kernel: no proof can unfold it.
\item \lead{Axioms and sorry} 0 declarations of these modules are axioms or use \code{sorryAx}.
\end{itemize}

\begin{tabular}{p{0.34\linewidth}p{0.14\linewidth}p{0.2\linewidth}p{0.22\linewidth}}
\toprule
Declaration & Kind & Result & Location \\
\midrule
\leandecl{instReprLBTerm.repr} & partial def & \leandecl{Std.Format} & \leanloc{LeanToLambdaBox/\allowbreak{}Basic.\allowbreak{}lean}{106} \\
\leandecl{LBTerm.inlinedSize} & partial def & \leandecl{Nat} & \leanloc{LeanToLambdaBox/\allowbreak{}Erasure.\allowbreak{}lean}{107} \\
\leandecl{LBTerm.isValue} & partial def & \leandecl{Bool} & \leanloc{LeanToLambdaBox/\allowbreak{}Erasure.\allowbreak{}lean}{123} \\
\leandecl{LBTerm.isValue.isConstructorApp} & partial def & \leandecl{Bool} & \leanloc{LeanToLambdaBox/\allowbreak{}Erasure.\allowbreak{}lean}{136} \\
\leandecl{LBTerm.stripLambdas} & partial def & \leandecl{LBTerm} & \leanloc{LeanToLambdaBox/\allowbreak{}Erasure.\allowbreak{}lean}{141} \\
\leandecl{Erasure.runMetaM} & def & \leandecl{Lean.Core.CoreM} & \leanloc{LeanToLambdaBox/\allowbreak{}Erasure.\allowbreak{}lean}{250} \\
\leandecl{Erasure.isErasable} & def & \leandecl{Lean.Meta.MetaM} & \leanloc{LeanToLambdaBox/\allowbreak{}Erasure.\allowbreak{}lean}{254} \\
\leandecl{Erasure.replaceUnsafeRecNames} & def & \leandecl{Lean.Core.CoreM} & \leanloc{LeanToLambdaBox/\allowbreak{}Erasure.\allowbreak{}lean}{270} \\
\leandecl{Erasure.prepareErasure} & def & \leandecl{Lean.Core.CoreM} & \leanloc{LeanToLambdaBox/\allowbreak{}Erasure.\allowbreak{}lean}{300} \\
\leandecl{Erasure.argMaskCore} & def & \leandecl{Lean.Core.CoreM} & \leanloc{LeanToLambdaBox/\allowbreak{}Erasure.\allowbreak{}lean}{323} \\
\leandecl{Erasure.erase} & def & \leandecl{Lean.Core.CoreM} & \leanloc{LeanToLambdaBox/\allowbreak{}Erasure.\allowbreak{}lean}{960} \\
\leandecl{Erasure.MLType.toString} & partial def & \leandecl{String} & \leanloc{LeanToLambdaBox/\allowbreak{}Erasure.\allowbreak{}lean}{979} \\
\leandecl{Erasure.MLType.toString.protArrow} & partial def & \leandecl{String} & \leanloc{LeanToLambdaBox/\allowbreak{}Erasure.\allowbreak{}lean}{989} \\
\leandecl{Erasure.MLType.toString.protCtor} & partial def & \leandecl{String} & \leanloc{LeanToLambdaBox/\allowbreak{}Erasure.\allowbreak{}lean}{992} \\
\leandecl{Erasure.to_ml_type} & partial def & \leandecl{Lean.Meta.MetaM} & \leanloc{LeanToLambdaBox/\allowbreak{}Erasure.\allowbreak{}lean}{998} \\
\leandecl{Erasure.gen_mli} & def & \leandecl{Lean.Meta.MetaM} & \leanloc{LeanToLambdaBox/\allowbreak{}Erasure.\allowbreak{}lean}{1014} \\
\leandecl{Erasure.eraseElab} & def & \leandecl{Lean.Elab.Command.CommandElab} & \leanloc{LeanToLambdaBox/\allowbreak{}Erasure/\allowbreak{}Command.\allowbreak{}lean}{18} \\
\leandecl{Erasure.eraseElab.unsafe_1} & unsafe def & \leandecl{Lean.Elab.Term.TermElabM} & \leanfile{LeanToLambdaBox/\allowbreak{}Erasure/\allowbreak{}Command.\allowbreak{}lean} \\
\leandecl{Erasure.eraseElab.unsafe_impl_2} & opaque, implemented\_\allowbreak{}by & \leandecl{Lean.Elab.Term.TermElabM} & \leanfile{LeanToLambdaBox/\allowbreak{}Erasure/\allowbreak{}Command.\allowbreak{}lean} \\
\leandecl{Erasure.eraseEntry} & def & \leandecl{Lean.Core.CoreM} & \leanloc{LeanToLambdaBox/\allowbreak{}Erasure/\allowbreak{}Entry.\allowbreak{}lean}{39} \\
\leandecl{LBTerm.to_sexpr} & partial def & \leandecl{sexpr} & \leanloc{LeanToLambdaBox/\allowbreak{}Printing.\allowbreak{}lean}{68} \\
\leandecl{edef.to_sexpr} & partial def & \leandecl{sexpr} & \leanloc{LeanToLambdaBox/\allowbreak{}Printing.\allowbreak{}lean}{87} \\
\leandecl{rocq_print} & def & \leandecl{IO} & \leanloc{LeanToLambdaBox/\allowbreak{}Printing.\allowbreak{}lean}{164} \\
\bottomrule
\end{tabular}



\section{Inherited trust: lean4lean}\label{sec:trust-lean4lean}

\begin{itemize}
\item \lead{Pin} the lean4lean package of Section~\ref{sec:intro-pins}.
\item \lead{Rule} A result built on lean4lean may depend on its sorry sources and axioms. Its node
  lists each of them (\emph{Inherited from lean4lean}); the audit checks the list against the Lean
  environment.
\end{itemize}

The tables are generated from the Lean environment by \code{blueprint/scripts/audit.py}.

% Generated by blueprint/scripts/audit.py --update from the Lean environment. Do not edit.
\begin{itemize}
\item \lead{Libraries} the modules Lean4Lean, Lean4Lean.\allowbreak{}Theory.\allowbreak{}*, Lean4Lean.\allowbreak{}Verify.\allowbreak{}* of the pinned lean4lean (22 sorry sources, 41 axioms).
\item \lead{Sorry source} a declaration whose own type or value uses \code{sorryAx}.
\item \lead{Reached} the blueprint nodes depend on \code{Lean4Lean.\allowbreak{}TrProj}, \code{Lean4Lean.\allowbreak{}VEnv.\allowbreak{}IsDefEqU.\allowbreak{}forallE\_\allowbreak{}inv\_\allowbreak{}stratified}, \code{Lean4Lean.\allowbreak{}VEnv.\allowbreak{}IsDefEqU.\allowbreak{}sort\_\allowbreak{}forallE\_\allowbreak{}inv}, \code{Lean4Lean.\allowbreak{}VEnv.\allowbreak{}IsDefEqU.\allowbreak{}sort\_\allowbreak{}inv}, \code{Lean4Lean.\allowbreak{}VEnv.\allowbreak{}addInduct}, \code{Lean4Lean.\allowbreak{}VEnv.\allowbreak{}addInduct\_\allowbreak{}WF}, \code{Lean4Lean.\allowbreak{}VInductDecl.\allowbreak{}WF}.
\end{itemize}

\begin{tabular}{p{0.45\linewidth}p{0.36\linewidth}l}
\toprule
Sorry source & Module & Kind \\
\midrule
\code{Lean4Lean.\allowbreak{}TrProj} & \code{Lean4Lean.\allowbreak{}Verify.\allowbreak{}Typing.\allowbreak{}Expr}:68 & def \\
\code{Lean4Lean.\allowbreak{}TrProj.\allowbreak{}defeqDFC} & \code{Lean4Lean.\allowbreak{}Verify.\allowbreak{}Typing.\allowbreak{}Lemmas}:725 & theorem \\
\code{Lean4Lean.\allowbreak{}TrProj.\allowbreak{}instL} & \code{Lean4Lean.\allowbreak{}Verify.\allowbreak{}Typing.\allowbreak{}Lemmas}:1508 & theorem \\
\code{Lean4Lean.\allowbreak{}TrProj.\allowbreak{}instN} & \code{Lean4Lean.\allowbreak{}Verify.\allowbreak{}Typing.\allowbreak{}Lemmas}:1239 & theorem \\
\code{Lean4Lean.\allowbreak{}TrProj.\allowbreak{}uniq} & \code{Lean4Lean.\allowbreak{}Verify.\allowbreak{}Typing.\allowbreak{}Lemmas}:936 & theorem \\
\code{Lean4Lean.\allowbreak{}TrProj.\allowbreak{}weak'} & \code{Lean4Lean.\allowbreak{}Verify.\allowbreak{}Typing.\allowbreak{}Lemmas}:641 & theorem \\
\code{Lean4Lean.\allowbreak{}TrProj.\allowbreak{}weak'\_\allowbreak{}inv} & \code{Lean4Lean.\allowbreak{}Verify.\allowbreak{}Typing.\allowbreak{}Lemmas}:722 & theorem \\
\code{Lean4Lean.\allowbreak{}TrProj.\allowbreak{}wf} & \code{Lean4Lean.\allowbreak{}Verify.\allowbreak{}Typing.\allowbreak{}Lemmas}:893 & theorem \\
\code{Lean4Lean.\allowbreak{}TypeChecker.\allowbreak{}Inner.\allowbreak{}inferProj.\allowbreak{}WF} & \code{Lean4Lean.\allowbreak{}Verify.\allowbreak{}TypeChecker.\allowbreak{}InferType}:388 & theorem \\
\code{Lean4Lean.\allowbreak{}TypeChecker.\allowbreak{}Inner.\allowbreak{}isDefEqUnitLike.\allowbreak{}WF} & \code{Lean4Lean.\allowbreak{}Verify.\allowbreak{}TypeChecker.\allowbreak{}IsDefEq}:486 & theorem \\
\code{Lean4Lean.\allowbreak{}TypeChecker.\allowbreak{}Inner.\allowbreak{}reduceProjCore.\allowbreak{}WF} & \code{Lean4Lean.\allowbreak{}Verify.\allowbreak{}TypeChecker.\allowbreak{}Reduce}:143 & theorem \\
\code{Lean4Lean.\allowbreak{}TypeChecker.\allowbreak{}Inner.\allowbreak{}reduceRecursor.\allowbreak{}WF} & \code{Lean4Lean.\allowbreak{}Verify.\allowbreak{}TypeChecker.\allowbreak{}WHNF}:6 & theorem \\
\code{Lean4Lean.\allowbreak{}TypeChecker.\allowbreak{}Inner.\allowbreak{}tryEtaStructCore.\allowbreak{}WF} & \code{Lean4Lean.\allowbreak{}Verify.\allowbreak{}TypeChecker.\allowbreak{}IsDefEq}:225 & theorem \\
\code{Lean4Lean.\allowbreak{}VEnv.\allowbreak{}IsDefEqU.\allowbreak{}forallE\_\allowbreak{}inv\_\allowbreak{}stratified} & \code{Lean4Lean.\allowbreak{}Theory.\allowbreak{}Typing.\allowbreak{}Injectivity}:14 & theorem \\
\code{Lean4Lean.\allowbreak{}VEnv.\allowbreak{}IsDefEqU.\allowbreak{}sort\_\allowbreak{}forallE\_\allowbreak{}inv} & \code{Lean4Lean.\allowbreak{}Theory.\allowbreak{}Typing.\allowbreak{}Injectivity}:33 & theorem \\
\code{Lean4Lean.\allowbreak{}VEnv.\allowbreak{}IsDefEqU.\allowbreak{}sort\_\allowbreak{}inv} & \code{Lean4Lean.\allowbreak{}Theory.\allowbreak{}Typing.\allowbreak{}Injectivity}:11 & theorem \\
\code{Lean4Lean.\allowbreak{}VEnv.\allowbreak{}IsDefEqU.\allowbreak{}weakN\_\allowbreak{}iff} & \code{Lean4Lean.\allowbreak{}Theory.\allowbreak{}Typing.\allowbreak{}UniqueTyping}:172 & theorem \\
\code{Lean4Lean.\allowbreak{}VEnv.\allowbreak{}NormalEq.\allowbreak{}parRed} & \code{Lean4Lean.\allowbreak{}Theory.\allowbreak{}Typing.\allowbreak{}ChurchRosser}:1178 & theorem \\
\code{Lean4Lean.\allowbreak{}VEnv.\allowbreak{}addInduct} & \code{Lean4Lean.\allowbreak{}Theory.\allowbreak{}Inductive}:7 & def \\
\code{Lean4Lean.\allowbreak{}VEnv.\allowbreak{}addInduct\_\allowbreak{}WF} & \code{Lean4Lean.\allowbreak{}Theory.\allowbreak{}Typing.\allowbreak{}InductiveLemmas}:8 & theorem \\
\code{Lean4Lean.\allowbreak{}VInductDecl.\allowbreak{}WF} & \code{Lean4Lean.\allowbreak{}Theory.\allowbreak{}Inductive}:5 & def \\
\code{Lean4Lean.\allowbreak{}addDecl.\allowbreak{}WF} & \code{Lean4Lean.\allowbreak{}Verify.\allowbreak{}Environment}:221 & theorem \\
\bottomrule
\end{tabular}

\begin{tabular}{p{0.3\linewidth}rp{0.55\linewidth}}
\toprule
Module & Axioms & Names \\
\midrule
\code{Lean4Lean.\allowbreak{}PtrEq} & 2 & \code{Lean4Lean.\allowbreak{}ptrEqConstantInfo\_\allowbreak{}eq}, \code{Lean4Lean.\allowbreak{}ptrEqExpr\_\allowbreak{}eq} \\
\code{Lean4Lean.\allowbreak{}Verify.\allowbreak{}Axioms} & 32 & \code{Lean.\allowbreak{}Expr.\allowbreak{}abstractRange\_\allowbreak{}eq}, \code{Lean.\allowbreak{}Expr.\allowbreak{}abstract\_\allowbreak{}eq}, \code{Lean.\allowbreak{}Expr.\allowbreak{}equal\_\allowbreak{}eq}, \code{Lean.\allowbreak{}Expr.\allowbreak{}eqv\_\allowbreak{}eq}, \code{Lean.\allowbreak{}Expr.\allowbreak{}hasLooseBVar\_\allowbreak{}eq}, \code{Lean.\allowbreak{}Expr.\allowbreak{}instantiate1\_\allowbreak{}eq}, \code{Lean.\allowbreak{}Expr.\allowbreak{}instantiateRange\_\allowbreak{}eq}, \code{Lean.\allowbreak{}Expr.\allowbreak{}instantiateRevRange\_\allowbreak{}eq}, \code{Lean.\allowbreak{}Expr.\allowbreak{}instantiateRev\_\allowbreak{}eq}, \code{Lean.\allowbreak{}Expr.\allowbreak{}instantiate\_\allowbreak{}eq}, \code{Lean.\allowbreak{}Expr.\allowbreak{}liftLooseBVars\_\allowbreak{}eq}, \code{Lean.\allowbreak{}Expr.\allowbreak{}looseBVarRange\_\allowbreak{}eq}, \code{Lean.\allowbreak{}Expr.\allowbreak{}lowerLooseBVars\_\allowbreak{}eq}, \code{Lean.\allowbreak{}Expr.\allowbreak{}mkAppData\_\allowbreak{}eq}, \code{Lean.\allowbreak{}Expr.\allowbreak{}mkData\_\allowbreak{}eq}, \code{Lean.\allowbreak{}Expr.\allowbreak{}replace\_\allowbreak{}eq}, \code{Lean.\allowbreak{}Level.\allowbreak{}hasMVar\_\allowbreak{}eq}, \code{Lean.\allowbreak{}Level.\allowbreak{}hasParam\_\allowbreak{}eq}, \code{Lean.\allowbreak{}Level.\allowbreak{}instLawfulBEqLevel}, \code{Lean.\allowbreak{}Level.\allowbreak{}isExplicitSubsumedAux\_\allowbreak{}eq}, \code{Lean.\allowbreak{}Level.\allowbreak{}mkData\_\allowbreak{}eq}, \code{Lean.\allowbreak{}Level.\allowbreak{}mkLevelIMaxCore\_\allowbreak{}eq}, \code{Lean.\allowbreak{}Level.\allowbreak{}mkMaxAux\_\allowbreak{}eq}, \code{Lean.\allowbreak{}Level.\allowbreak{}normalize\_\allowbreak{}eq}, \code{Lean.\allowbreak{}Level.\allowbreak{}skipExplicit\_\allowbreak{}eq}, \code{Lean.\allowbreak{}PersistentArray.\allowbreak{}toList'\_\allowbreak{}push}, \code{Lean.\allowbreak{}PersistentHashMap.\allowbreak{}WF.\allowbreak{}find?\_\allowbreak{}eq}, \code{Lean.\allowbreak{}PersistentHashMap.\allowbreak{}WF.\allowbreak{}toList'\_\allowbreak{}insert}, \code{Lean.\allowbreak{}PersistentHashMap.\allowbreak{}findAux\_\allowbreak{}isSome}, \code{Lean.\allowbreak{}Syntax.\allowbreak{}structEq\_\allowbreak{}eq}, \code{Std.\allowbreak{}TreeMap.\allowbreak{}all\_\allowbreak{}eq\_\allowbreak{}all\_\allowbreak{}toList}, \code{Std.\allowbreak{}TreeMap.\allowbreak{}any\_\allowbreak{}eq\_\allowbreak{}any\_\allowbreak{}toList} \\
\code{Lean4Lean.\allowbreak{}Verify.\allowbreak{}Expr} & 4 & \code{Lean.\allowbreak{}Expr.\allowbreak{}Data.\allowbreak{}looseBVarRange\_\allowbreak{}le.\allowbreak{}\_\allowbreak{}native.\allowbreak{}bv\_\allowbreak{}decide.\allowbreak{}ax\_\allowbreak{}1\_\allowbreak{}7}, \code{Lean.\allowbreak{}Expr.\allowbreak{}mkAppData\_\allowbreak{}looseBVarRange.\allowbreak{}\_\allowbreak{}native.\allowbreak{}bv\_\allowbreak{}decide.\allowbreak{}ax\_\allowbreak{}1\_\allowbreak{}8}, \code{Lean.\allowbreak{}Expr.\allowbreak{}mkData\_\allowbreak{}looseBVarRange.\allowbreak{}\_\allowbreak{}native.\allowbreak{}bv\_\allowbreak{}decide.\allowbreak{}ax\_\allowbreak{}1\_\allowbreak{}9}, \code{\_\allowbreak{}private.\allowbreak{}Lean4Lean.\allowbreak{}Verify.\allowbreak{}Expr.\allowbreak{}0.\allowbreak{}Lean.\allowbreak{}Expr.\allowbreak{}mkData\_\allowbreak{}flags.\allowbreak{}\_\allowbreak{}native.\allowbreak{}bv\_\allowbreak{}decide.\allowbreak{}ax\_\allowbreak{}1\_\allowbreak{}12} \\
\code{Lean4Lean.\allowbreak{}Verify.\allowbreak{}Level} & 3 & \code{Lean.\allowbreak{}Level.\allowbreak{}mkData\_\allowbreak{}depth.\allowbreak{}\_\allowbreak{}native.\allowbreak{}bv\_\allowbreak{}decide.\allowbreak{}ax\_\allowbreak{}1\_\allowbreak{}9}, \code{Lean.\allowbreak{}Level.\allowbreak{}mkData\_\allowbreak{}hasMVar.\allowbreak{}\_\allowbreak{}native.\allowbreak{}bv\_\allowbreak{}decide.\allowbreak{}ax\_\allowbreak{}1\_\allowbreak{}8}, \code{Lean.\allowbreak{}Level.\allowbreak{}mkData\_\allowbreak{}hasParam.\allowbreak{}\_\allowbreak{}native.\allowbreak{}bv\_\allowbreak{}decide.\allowbreak{}ax\_\allowbreak{}1\_\allowbreak{}8} \\
\bottomrule
\end{tabular}



\section*{Caveats}
\begin{itemize}
\item \lead{Caveat} A green shipping node means that its declarations exist and use standard axioms
  only. It says nothing about their behaviour.
\item \lead{Caveat} The eraser runs as Lean meta-code in the elaborator. Its execution, and Lean's
  compiler passes that it calls, are outside every Lean statement.
\item \lead{Caveat} A \code{panic!} or \code{unreachable!} of the eraser returns a default value and
  continues (register entries R-4 and R-5, Chapter~\ref{chap:shipping-changes}).
\end{itemize}
